% !TEX root = ../Thesis.tex
% ============================================================
\chapter{Results} \label{chapter:results}
% ============================================================

% Do not add a number here without a command that reproduces it.

\section{What Is Measured Here, and What Is Not} \label{sec:results-scope}

The results below come from two sources, and conflating them would make the
chapter worthless.

\paragraph{Measurements on the clinical graph.} Sections~\ref{sec:results-graph}
and~\ref{sec:results-preparation} report properties of the live STCS graph and of
one patient loaded from it, obtained through read-only queries on hospital
infrastructure. They are facts about the data.

\paragraph{Experiments on a shape-matched replica.}
Sections~\ref{sec:results-dilution} to~\ref{sec:results-candidates} report
experiments run on a synthetic patient generated to reproduce the measured
structure of one real patient exactly, with the same node and edge counts, the same
label distribution, the same ratio of nodes to distinct texts, the same
proportion of diagnoses lacking a description. The generator emits rows in the
shape the database returns and passes them through the same assembly code as the
real loader, so the structure is identical by construction rather than by
imitation. Its \emph{content}, however, is invented: analyte names, drug names
and diagnosis labels are drawn from synthetic vocabularies.

The consequence is a strict limit on what these experiments can support. They are
evidence about \emph{mechanism}, meaning whether a size parameter controls size,
whether tie-breaking decides an outcome, whether resolving a code changes what can
be retrieved, since those depend on structure. They are not evidence about
clinical retrieval quality, which depends on content, and no recall figure from
the replica should be read as an estimate of recall on the cohort.

\paragraph{Not yet measured.} The comparison against SnapQuery, which
Chapter~\ref{chapter:introduction} states as RQ1, is not reported here. Neither
is any result using a trained sentence encoder: every experiment below uses a
deterministic bag-of-words encoder, which is adequate for the structural
questions asked and inadequate for anything about semantic quality. Both are
addressed in Section~\ref{sec:results-threats}.

\section{The Graph as Measured} \label{sec:results-graph}

Table~\ref{tab:measured-graph} collects the properties of the cohort graph that
drive every design decision in this thesis.

\begin{table}[t]
\centering
\caption{The STCS cohort graph at University Hospital Basel, measured through
read-only queries. Percentages are of the whole graph.}
\label{tab:measured-graph}
\begin{tabular}{lr}
\toprule
Property & Value \\
\midrule
Nodes & 21{,}537{,}305 \\
Edges & 72{,}218{,}890 \\
Patients & 1{,}197 \\
Nodes per patient (mean) & $\approx 17{,}993$ \\
\midrule
\multicolumn{2}{l}{\emph{Concentration of the graph in one construct}} \\
\texttt{Sample} + \texttt{LabTest} + \texttt{LabResult} + \texttt{LabTestEvent} & 20{,}523{,}891 (95.3\%) \\
\midrule
\multicolumn{2}{l}{\emph{Degenerate connectivity}} \\
\texttt{SourceSystem} nodes / mean degree & 168 / 69{,}525 \\
\texttt{SourceSystem} maximum degree & 10{,}858{,}404 \\
\texttt{BodySite} nodes / edges carried & 6 / 3{,}454{,}881 \\
\texttt{TimePattern} nodes / mean degree & 2 / 313{,}012 \\
\midrule
\multicolumn{2}{l}{\emph{Relation multiplicity}} \\
\texttt{hasSourceSystem} & 11{,}680{,}076 \\
\texttt{hasSubjectPseudoIdentifier} & 11{,}600{,}854 \\
\texttt{hasAdministrativeCase} & 11{,}566{,}354 \\
\texttt{hasCode} & 4{,}883{,}809 \\
\bottomrule
\end{tabular}
\end{table}

Two entries deserve comment. A single \texttt{SourceSystem} node carries
$10{,}858{,}404$ edges, fifteen per cent of every edge in the graph, and the four
relations listed account for roughly half of all edges between them. This is not
a data quality problem: provenance is required by the SPHN schema and is
correctly recorded. It is, however, fatal to an unqualified connectivity
requirement, and it is the empirical basis for Prediction~3 of
Section~\ref{sec:pcst-consequences}.

\paragraph{Terminology coverage.} Table~\ref{tab:terminology-coverage} shows the
five terminologies present and how much of each carries a description in the
graph. The asymmetry is the important part: laboratory codes are fully described
and diagnosis codes are two-thirds undescribed.

\begin{table}[t]
\centering
\caption{Terminologies in the cohort graph and their description coverage.
ICD-10-GM is the German national modification of ICD-10 and its labels are in
German; LOINC descriptions are in English.}
\label{tab:terminology-coverage}
\begin{tabular}{llrr}
\toprule
Terminology & Used for & Distinct codes & With a description \\
\midrule
LOINC~\cite{loinc} & laboratory tests & 544 & 544 (100\%) \\
ICD-10-GM~\cite{icd10} & diagnoses & 2{,}921 & 972 (33\%) \\
ATC~\cite{atc} & drugs & 737 & not measured \\
UCUM & units & 70 & not measured \\
SNOMED~CT~\cite{snomedct} & mixed & 69 & not measured \\
\bottomrule
\end{tabular}
\end{table}

\section{Graph Preparation} \label{sec:results-preparation}

Table~\ref{tab:patient-zero} reports one patient after the preparation procedure
of Chapter~4: shared terminology nodes folded into the text of the events that
reference them, provenance pruned, and the four-node laboratory construct
collapsed to one node per observation.

\begin{table}[t]
\centering
\caption{Patient \#0 after graph preparation. The patient index is a position in
a query result, not an identifier.}
\label{tab:patient-zero}
\begin{tabular}{lr}
\toprule
Quantity & Value \\
\midrule
Nodes & 15{,}810 \\
Edges & 45{,}562 \\
\midrule
\texttt{LabObservation} & 12{,}349 \\
\texttt{LabTestType} (distinct analytes) & 198 \\
\texttt{DrugAdministrationEvent} & 2{,}655 \\
\texttt{Drug} (distinct) & 279 \\
\texttt{BilledDiagnosis} & 222 \\
\texttt{AdministrativeCase} & 106 \\
\midrule
Distinct \texttt{embed\_text} values & 614 \\
Nodes per distinct text & 25.7 \\
Diagnoses with no description & 178 of 222 (80\%) \\
Laboratory tests with no description & 0 of 12{,}349 \\
\bottomrule
\end{tabular}
\end{table}

Three observations follow.

First, preparation is a substantial reduction on the dominant construct. SPHN
represents one laboratory result as four event nodes plus shared quantity, unit,
code and reference-range nodes; the loader emits one. For this patient that turns
roughly $49{,}400$ raw laboratory event nodes into $12{,}349$.

Second, one prepared patient is $15{,}810$ nodes, about eleven times the average
WebQSP sample of $1{,}371$ nodes on which G-Retriever was
validated~\cite{he2024gretriever}, and this is one patient of $1{,}197$. The
narrowing step of Section~\ref{sec:results-candidates} is a consequence of this
row, not a convenience.

Third, $15{,}810$ nodes carry $614$ distinct descriptions. A similarity score can
therefore sort this graph into at most $614$ classes, and within a class the
ordering is decided by the sorting routine. This is Prediction~2 of
Section~\ref{sec:pcst-consequences}, measured.

\section{Prediction 1: The Edge-Prize Mechanism Collapses} \label{sec:results-dilution}

Section~\ref{sec:pcst-consequences} predicted from
Equations~\ref{eq:edge-prizes} and~\ref{eq:cost-adjustment} that heavily shared
relation types drive the effective edge cost towards zero, so that PCST returns
its entire input. Running the published default configuration
($\texttt{topk}=3$, $\texttt{topk\_e}=5$, $c_e = 0.5$) on the shape-matched
replica of patient \#0 returned \textbf{15{,}810 of 15{,}810 nodes and
45{,}562 of 45{,}562 edges}, which is the whole graph, for every question tested.

The prediction is confirmed, and the remedy follows directly from the algebra
rather than from tuning: setting $\texttt{topk\_e} = 0$ removes edge prizes
entirely, which leaves $c_e$ free of the cap in
Equation~\ref{eq:cost-adjustment}. Under that configuration the same solver on
the same graph returns eleven nodes. Every subsequent experiment in this chapter
uses $\texttt{topk\_e} = 0$ with an explicit $c_e$ for this reason, and the
choice is an ablation switch rather than a hidden default.

\section{Size Control Across Selection Families} \label{sec:results-size}

Table~\ref{tab:sweep} sweeps each family across its own size parameter on the
full-size replica and reports the subgraph actually produced. Because the
parameters are not the same quantity, being a cost against a count against a
hop number, the comparison is made on measured size, as
Section~\ref{sec:methodology-comparability} requires.

\begin{table}[t]
\centering
\caption{Recall against measured size on the shape-matched replica, 14 planted
questions. ``Fraction'' is of the 15{,}810-node graph. Recall figures compare
only at comparable size; the point of the table is the size column.}
\label{tab:sweep}
\begin{tabular}{llrrrc}
\toprule
Family & Dial & Nodes & Fraction & Recall & Connected \\
\midrule
PCST & $c_e = 3.0$   & 8.0  & 0.05\% & 0.158 & yes \\
PCST & $c_e = 1.0$   & 10.0 & 0.06\% & 0.162 & yes \\
PCST & $c_e = 0.5$   & 11.0 & 0.07\% & 0.164 & yes \\
PCST & $c_e = 0.1$   & 11.0 & 0.07\% & 0.164 & yes \\
PCST & $c_e = 0.01$  & 11.0 & 0.07\% & 0.164 & yes \\
PCST & published defaults & 15{,}810.0 & 100\% & 1.000 & yes \\
\midrule
Shortest paths & $k = 2$  & 3.0  & 0.02\% & 0.147 & yes \\
Shortest paths & $k = 10$ & 11.0 & 0.07\% & 0.164 & yes \\
\midrule
Top-$k$ nodes + neighbours & $k = 10$  & 21.7    & 0.14\%  & 0.164 & yes \\
Top-$k$ nodes + neighbours & $k = 30$  & 2{,}243.7 & 14.19\% & 0.277 & yes \\
Top-$k$ nodes + neighbours & $k = 100$ & 4{,}514.9 & 28.56\% & 0.503 & yes \\
\midrule
BFS expansion & $h = 1$ & 7.7        & 0.05\%  & 0.150 & yes \\
BFS expansion & $h = 2$ & 15{,}333.7 & 96.99\% & 0.925 & yes \\
BFS expansion & $h = 3$ & 15{,}810.0 & 100\%   & 1.000 & yes \\
\midrule
Top-$k$ triples (KAPING) & $k = 10$  & 11.0  & 0.07\% & 0.164 & yes \\
Top-$k$ triples (KAPING) & $k = 100$ & 93.6  & 0.59\% & 0.352 & \textbf{no} \\
Top-$k$ triples (KAPING) & $k = 300$ & 209.9 & 1.33\% & 0.519 & \textbf{no} \\
\bottomrule
\end{tabular}
\end{table}

\begin{figure}[t]
\centering
\includegraphics[width=\textwidth]{recall-vs-size}
\caption[Answer-node recall against measured subgraph size]{Answer-node recall
against measured subgraph size, one line per selection strategy, each swept
across its own parameter range. Size is plotted logarithmically because the
configurations span three orders of magnitude, and against the size actually
achieved rather than against each strategy's own parameter, which are not
commensurable (Section~\ref{sec:methodology-comparability}). Reading a vertical
slice gives the comparison the thesis asks for: at a fixed budget, which
strategy recalls most? Shortest paths appears only at the far left because it
cannot be made to return a large subgraph. Generated by
\texttt{experiments/figures.py}.}
\label{fig:recall-vs-size}
\end{figure}

Four results, in decreasing order of how surprising they are.

\paragraph{At equal size, connectivity contributed nothing.} At approximately
eleven nodes, PCST, shortest paths, top-$k$ nodes with neighbours and KAPING all
scored $0.164$, identically, and Figure~\ref{fig:recall-vs-size} marks the point
where the four curves meet. On this graph, at the sizes PCST operates in, what
gets retrieved is decided by the similarity signal and not by the structural
strategy. This is a direct partial answer to RQ2, and it is a negative one.

\paragraph{Neighbourhood expansion has no usable middle.} BFS returns $7.7$ nodes
at one hop and $15{,}334$ at two, rising from $0.05\%$ of the graph to $97\%$ with no
setting in between, because the dial is an integer. Top-$k$ nodes with
neighbours behaves the same way once $k$ is large enough to catch a hub: $21.7$
nodes at $k=10$, $2{,}244$ at $k=30$. The claim made in
Section~\ref{sec:subgraph-selection-approaches}, that these methods lose size
control on a dense graph, is here measured rather than asserted.

\paragraph{PCST's size dial is narrower than advertised.} With
$\texttt{topk} = 10$, sweeping $c_e$ from $3.0$ down to $0.01$ moved the subgraph
from $8$ nodes to $11$ and then stopped. The reachable range is $8$--$11$ nodes:
$c_e$ saturates in both directions, upward by Equation~\ref{eq:cost-adjustment}
and downward because once every prized node is affordable there is nothing
further to buy. The effective size dial in this regime is $\texttt{topk}$, not
$c_e$, and a methodology that sweeps only $c_e$, as the published description
suggests, explores a narrow band and reports it as a curve.
Figure~\ref{fig:pcst-size-dial} shows the whole of that band.

\begin{figure}[t]
\centering
\includegraphics[width=0.82\textwidth]{pcst-size-dial}
\caption[Subgraph size against the PCST edge cost]{Subgraph size against PCST's
edge cost $c_e$ across the published range, holding \texttt{topk} fixed. The
entire sweep moves the result by three nodes. A parameter presented as the size
dial does not, in this regime, control size.}
\label{fig:pcst-size-dial}
\end{figure}

\paragraph{Independent scoring returns disconnected results.} KAPING's output is
disconnected for every $k \geq 100$, as the family's description predicts. It is
nonetheless competitive on recall at matched size, which is worth stating plainly
rather than burying: the connectivity guarantee is not, on this graph, buying
recall.

\section{Terminology Resolution} \label{sec:results-terminology}

Section~\ref{subsec:terminologies} argued that resolving codes to descriptions is
a precondition for retrieval rather than preprocessing hygiene.
Table~\ref{tab:terminology-gain} measures it, using questions that never quote
the code they are looking for, since a question containing \texttt{N18.5} finds the
\texttt{N18.5} node by string overlap and measures nothing.

\begin{table}[t]
\centering
\caption{Answer-node recall on diagnosis questions under a catalogue covering
85\% of code-only codes across all resolution tiers. Identical for both
retrievers, which is itself the result.}
\label{tab:terminology-gain}
\begin{tabular}{lcc}
\toprule
Condition & Described diagnoses & Code-only diagnoses \\
\midrule
No catalogue & 1.00 & \textbf{0.00} \\
Catalogue, language matches the question & 1.00 & \textbf{0.80} \\
Catalogue, plus parent-concept words & 1.00 & 0.80 \\
Catalogue, language differs from the question & \textbf{0.00} & \textbf{0.00} \\
\bottomrule
\end{tabular}
\end{table}

\begin{figure}[t]
\centering
\includegraphics[width=\textwidth]{terminology-gain}
\caption[Terminology resolution by catalogue condition]{Answer-node recall by
catalogue condition, split by whether the diagnosis already carried a name.
The split is the point: averaging the two halves would hide each effect behind
the other. Values are printed at the end of every bar because a recall of
exactly $0.00$ draws no bar, and here those zeros are measurements rather than
missing conditions.}
\label{fig:terminology-gain}
\end{figure}

Three conclusions.

Resolution moves code-only diagnoses from unreachable to $0.80$; the shortfall
from $1.00$ is the deliberately incomplete catalogue, and the per-tier counts are
reported by the enrichment step rather than folded into a single coverage
figure.

The gain is \emph{identical} for PCST and for KAPING. Terminology resolution
changes what is reachable at all, not which method reaches it better, so it
belongs to graph preparation and must not be reported as a retrieval result.

A catalogue in the wrong language is worse than no catalogue: it replaced text
that was already matching and took described diagnoses from $1.00$ to $0.00$.
Whether this step helps or harms therefore depends on a fact about deployment,
namely the language users write their questions in, and that is not a detail to
settle late.

Adding parent-concept words, the simplest form of the schema signal of RQ3,
changed nothing here. The synthetic vocabulary gives leaf labels that already
subsume their parents' words, so the replica cannot fairly judge this; it needs a
catalogue with real hierarchical labels.

\section{Candidate Generation} \label{sec:results-candidates}

One prepared patient is far larger than the per-question subgraphs G-Retriever
was built for, so Chapter~4 introduces a narrowing stage.
Table~\ref{tab:candidates} reports its effect. \emph{Ceiling} is the recall the
inner retriever could have reached given what the narrowing kept: reporting
recall without it makes a selection failure and a narrowing failure look
identical.

\begin{table}[t]
\centering
\caption{Two-stage retrieval on the full-size replica, PCST inner. Recall is
comparable only at comparable node count, which is why that column is present.}
\label{tab:candidates}
\begin{tabular}{lrrrr}
\toprule
Stage one & Nodes & Recall & Ceiling & ms / question \\
\midrule
none (whole patient graph)  & 11.0  & 0.164 & 1.000 & 78 \\
Seed expansion, cap 500     & 11.0  & 0.164 & 0.643 & 19 \\
Seed expansion, cap 2{,}000 & 11.0  & 0.164 & 0.773 & 24 \\
Top-$n$ similar, $n = 500$  & 220.6 & 0.382 & 0.648 & 14 \\
\bottomrule
\end{tabular}
\end{table}

\begin{figure}[t]
\centering
\includegraphics[width=\textwidth]{candidate-generation}
\caption[Recall beside its ceiling for each narrowing condition]{Recall beside
the ceiling imposed by candidate generation, averaged over the question set. The
ceiling is the recall attainable given what stage one kept, so the gap between
the pair is the selector's shortfall and the ceiling itself is the generator's.
Reporting recall without it would make a narrowing failure indistinguishable
from a selection failure, and the two need opposite fixes.}
\label{fig:candidate-generation}
\end{figure}

Narrowing by seed expansion is free: the same eleven-node subgraph at the same
recall, three to four times faster. Across $1{,}197$ patients that is the
difference between a parameter sweep that completes and one that does not.

Top-$n$ similarity does not control size. The most similar nodes are not adjacent
to one another, so the induced region is nearly edgeless, PCST is handed isolated
fragments and returns all of them: $220.6$ nodes rather than $11$. Its higher
recall is a larger subgraph, not a better one. This is the clearest illustration in
this thesis of why recall must always be read beside a size.

The finding that was not anticipated is in the ceiling column when broken out by
question type: \textbf{every generator discarded every code-only diagnosis},
ceiling $0.000$. A node with no matching text is neither similar to the question
nor adjacent to anything that is, so narrowing removes it first. Terminology
resolution is therefore a \emph{prerequisite} for candidate generation, not an
independent improvement alongside it, and the two steps must be ordered.

\section{Threats to Validity} \label{sec:results-threats}

\paragraph{The comparison this thesis exists to make is not in this chapter.}
RQ1 asks whether PCST-based selection outperforms SnapQuery, and no SnapQuery
result is reported. Until one is, the chapter establishes how the method behaves
on this data and not whether it is preferable to what is deployed.

\paragraph{The encoder is a placeholder.} Every experiment uses a deterministic
bag-of-words encoder. That is appropriate for structural questions, since a size
dial either controls size or does not, whatever the embedding, and inappropriate
for
anything semantic. In particular the cross-language result of
Section~\ref{sec:results-terminology} shows that a \emph{lexical} encoder cannot
bridge a language gap; a multilingual sentence encoder might, and that experiment
has not been run.

\paragraph{Content is synthetic.} The replica is faithful in structure and
invented in content, so the recall values are properties of a synthetic
vocabulary. Comparisons between rows of the same table are meaningful; the
absolute values are not estimates of clinical performance.

\paragraph{One patient.} Every experiment derives from the structure of a single
patient, who may not be typical of the cohort. The distribution of patient sizes
across the $1{,}197$ has not been measured, and it determines whether one graph
per question is feasible at all.

\paragraph{Questions are planted, not collected.} The evaluation questions are
generated with known answers rather than drawn from real usage. They test the
harness and the mechanism; they cannot speak to what clinicians actually ask, and
in particular the finding that a question such as ``all creatinine values'' has
$2{,}407$ answer nodes for one patient suggests that a share of real questions
are aggregations, for which answer-node recall is undefined
(Section~\ref{sec:thesis-scope}).
